% SPDX-License-Identifier: CC-BY-4.0
% Copyright (C) 2026 Fabian Januszewski, Abhijit (the Langlandscape team)
% Frame 3: what we want: an integral Hecke-equivariant map, (a)-(e); scope.
% Sources: submission brief (FABIAN'S QUESTION, scope addendum); final report, introduction.
\begin{frame}{What we want: an integral Hecke-equivariant \emph{map}}
  \small
  {\footnotesize Desired outcome: ``A new approach to the construction of the Jacquet-Langlands functorial transfer and the characterization of its image.''}\\[2pt]
  A simplicially realisable, computable transfer by Hecke-equivariant maps, defined \emph{integrally},
  \[H^0(\text{quaternionic})_\m\ \longleftrightarrow\ H^1(\GL_2)_\m\qquad(\m\ \text{non-Eisenstein}),\]
  such that
  \begin{enumerate}\setlength{\itemsep}{0pt}
    \item[(a)] $\otimes\,\mathbb C$ gives an explicit instance of cohomological JL;
    \item[(b)] over $\ZZ_p$, after $\otimes\,\QQ_p$, gives Chenevier's $p$-adic JL;
    \item[(c)] over $\ZZ/p^k$, on ordinary parts, gives $\Lambda$-adic JL;
    \item[(d)] in (c), congruence modules can be defined and described;
    \item[(e)] everything is explicitly computable (in the spirit of HACS).
  \end{enumerate}
  \textbf{Why a map.} A trace identity matches spectra. A map transports lattices, hence congruences, and its \emph{kernel and
  cokernel} carry the congruence information that the identity cannot see.\\[3pt]
  \textbf{Scope.} JL holds for every inner form of $\GL_2$ over every number field (and for $\GL_n$);\\
  we work over $\QQ$ with
  \emph{definite} quaternion algebras only.
\end{frame}
